%% ============================================================
%% appendix_str.tex — 부록 F: 소프트웨어 테스트 결과 보고서 (STR v1.1)
%% 원문: Test/3_STR.html (2026-07-08, 1차 실측 결과 취합 반영판)
%% TCL v1.0 (14 카테고리 · 98 케이스) — 실기(BLE) / Web·Sim_Parts 시뮬레이터 병행
%% ============================================================
\chaptersummary{TCL 98개 케이스를 1차로 실측한 소프트웨어 테스트 결과 보고서 STR v1.1이다. 85건 Pass·0건 Fail 결과와 조건부 Go 판정을 담는다.}
\chapter{소프트웨어 테스트 결과 보고서 (STR v1.1)}
\label{app:str}

본 부록은 11일차에 수립한 QA 체계$\cdot$통합 테스트 케이스 명세서(TCL v1.0, 14 카테고리
$\cdot$98 케이스)에 따라 \textbf{12일차에 본격 수행한 1차 소프트웨어 테스트의 결과 보고서
(STR v1.1)}를 정리한 것이다. 테스트는 하드웨어 의존 항목은 \textbf{실기(BLE)}로, 그 외는
\textbf{Web/Sim\_Parts 시뮬레이터}로 병행 수행했으며, 테스터가 TCL 실행형 시트에
기록$\cdot$내보내기한 실측 결과와 미결 13건의 코멘트(r.txt)를 취합해 작성되었다.
원문은 \code{Test/3\_STR.html}.

\section{핵심 요약}
TCL v1.0 98 케이스에 대한 1차 실측 결과, \textbf{85건 Pass $\cdot$ 0건 Fail $\cdot$
7건 보류 $\cdot$ 3건 기타 $\cdot$ 3건 미수행}을 기록했다(전체 Pass 비율 \textbf{86.7\%}).
실행 완료된 85건 기준 \textbf{Fail 0건}으로 기능 결함은 발견되지 않았다. 미결 13건의
원인은 제품 결함이 아니라 \textbf{테스트 인프라(Web 수신 로그 미표시) 5건 $\cdot$ 장비
부족(총 모듈, 다수 BLE 기기) 4건 $\cdot$ 명세/절차 문제 2건 $\cdot$ 레거시/미노출 코드
2건}으로 분석되었다. High 우선순위 케이스 5건이 보류 상태이므로 판정은 \textbf{조건부
Go} --- Web 수신(rx) 로그 표시 기능 반영과 보류 재시험 완료를 조건으로 최종 확정한다.

\section{테스트 환경}
\begin{table}[htbp]\centering\small
\caption{테스트 환경}
\renewcommand{\arraystretch}{1.2}
\begin{tabularx}{\linewidth}{@{}l X@{}}
\toprule
\sffamily 구분 & \sffamily 상세 \\
\midrule
OS / 브라우저 & Windows 10/11, Chrome$\cdot$Edge 최신(Web Bluetooth 지원) \\
하드웨어 & Raspberry Pi Pico, H-브릿지 dual PWM, HC-SR04, 연속회전 서보$\times$2,
  LED 6개, 부저, 디지털 자기센서, OLED SSD1306, HM-10/BT05 \\
통신 & BLE 4.2+ / UART 9600 baud \\
시뮬레이션 & Web/Sim\_Parts (\code{simulateWorkspace}, 4개 주제) \\
\bottomrule
\end{tabularx}
\end{table}

이번 회차에서 확인된 환경 제약(미결 발생 원인): ① \textbf{Web 수신(rx) 로그 미표시} ---
로그가 송신 명령 위주라 Pico 회신 문자열(\code{1}, \code{READY})을 화면에서 확인할 수 없어
``응답 수신 확인''이 합격 기준인 케이스 5건 보류, ② \textbf{총(발사) 모듈 미확보} ---
전투 3건 미수행, ③ \textbf{BLE 테스트 베드 부족} --- 다수$\cdot$이종 기기 필터 검증 불가,
④ \textbf{계측 장비(오실로스코프) 미투입} --- 성능 KPI 정밀 실측 이월.

\section{수행 결과 요약 · 카테고리별 결과}
\begin{table}[htbp]\centering\small
\caption{카테고리별 결과 (총 98 케이스)}
\renewcommand{\arraystretch}{1.18}
\begin{tabularx}{\linewidth}{@{}X c c c c c c r@{}}
\toprule
\sffamily 카테고리 & \sffamily TC & \sffamily Pass & \sffamily Fail & \sffamily 보류 & \sffamily 기타 & \sffamily 미수행 & \sffamily Pass 비율 \\
\midrule
연결$\cdot$시스템        & 7  & 5  & 0 & 2 & 0 & 0 & 71.4\% \\
BLE$\cdot$파이프라인     & 8  & 5  & 0 & 3 & 0 & 0 & 62.5\% \\
BATCH 블록              & 6  & 5  & 0 & 1 & 0 & 0 & 83.3\% \\
블록 코딩 엔진           & 7  & 6  & 0 & 0 & 1 & 0 & 85.7\% \\
서보 모터 휠             & 6  & 6  & 0 & 0 & 0 & 0 & \textbf{100\%} \\
DC 모터                 & 5  & 5  & 0 & 0 & 0 & 0 & \textbf{100\%} \\
LED 6개                 & 6  & 5  & 0 & 0 & 1 & 0 & 83.3\% \\
디스플레이 OLED          & 6  & 6  & 0 & 0 & 0 & 0 & \textbf{100\%} \\
소리$\cdot$부저          & 5  & 4  & 0 & 1 & 0 & 0 & 80.0\% \\
전투(총)                & 3  & 0  & 0 & 0 & 0 & 3 & 미수행 \\
센서 입력               & 5  & 5  & 0 & 0 & 0 & 0 & \textbf{100\%} \\
설정$\cdot$안정성        & 6  & 5  & 0 & 0 & 1 & 0 & 83.3\% \\
시뮬레이터 3D           & 24 & 24 & 0 & 0 & 0 & 0 & \textbf{100\%} \\
시뮬레이터$\cdot$프리셋 E2E & 4 & 4 & 0 & 0 & 0 & 0 & \textbf{100\%} \\
\midrule
\sffamily 합계 & \textbf{98} & \textbf{85} & \textbf{0} & \textbf{7} & \textbf{3} & \textbf{3} & \textbf{86.7\%} \\
\bottomrule
\end{tabularx}
\end{table}

서보$\cdot$DC$\cdot$OLED$\cdot$센서$\cdot$\textbf{시뮬레이터 3D(24건 전체)}$\cdot$E2E 등
하드웨어 동작과 시뮬레이터 경로는 \textbf{전건 Pass} --- 핵심 교육 시나리오(주행$\cdot$표시$\cdot$%
측정$\cdot$시뮬레이션)의 완성도가 실측으로 확인되었다. 보류가 집중된 연결$\cdot$BLE
카테고리는 기능 이상이 아니라 \textbf{응답 문자열 확인 수단 부재}와 테스트 베드 부족이
원인이다. 전투 카테고리는 총 모듈 수급 후 일괄 재시험 예정이며, \code{gun=None} 상태의
부팅 안정성(TC-SYS-005)은 Pass로 확인되어 모듈 부재가 다른 기능에 주는 영향은 없다.

\section{미결 케이스 분석 (13건)}
\begin{table}[htbp]\centering\small
\caption{미결 13건의 원인 유형 분류}
\renewcommand{\arraystretch}{1.2}
\begin{tabularx}{\linewidth}{@{}l X c X@{}}
\toprule
\sffamily 유형 & \sffamily 정의 & \sffamily 건수 & \sffamily 해당 케이스 \\
\midrule
A. 테스트 인프라 & Web UI가 수신 응답을 표시하지 않아 합격 기준 검증 불가 & 5 &
  TC-CONN-001$\cdot$002, TC-BLE-002, TC-BAT-002, TC-SND-003 \\
B. 장비$\cdot$환경 & 필요한 하드웨어$\cdot$기기 구성을 확보하지 못함 & 4 &
  TC-BLE-001, TC-GUN-001$\sim$003 \\
C. 명세$\cdot$절차 & TCL 절차 안내$\cdot$전제조건 기술이 불충분 & 2 &
  TC-BLE-006, TC-BLK-005 \\
D. 레거시$\cdot$미노출 & 명령이 펌웨어에 있으나 Web 실행 경로(블록$\cdot$UI)가 없음 & 2 &
  TC-LED-005, TC-SYS-004 \\
\bottomrule
\end{tabularx}
\end{table}

\textbf{유형 A(보류 5건).} 5건 모두 테스터가 ``작동으로 보이지만 현재 Web 로그 출력
방식상 출력을 확인할 수 없습니다''로 동일하게 기록 --- 물리 동작(OLED 표시, 순차 실행,
멜로디 재생)은 관찰되나 합격 기준의 응답 문자열(\code{1}, \code{READY}) 대조가 불가하다.
Web 로그가 송신(tx) 위주이고 \code{bluetooth.js} notification 수신부의 rx 문자열을
노출하지 않는 것이 원인으로, 시뮬레이터 경로가 →송신/↩응답을 모두 표기(TC-SIM-004
Pass)하는 것과 대조적이다. 응답 대기형 명령이 \code{RESPONSE\_TIMEOUT}(5000ms) 없이
진행된 점은 응답 정상 수신의 간접 증거다. 조치: 실기 BLE 경로에도 수신 문자열 로그
표시(IMP-001, High) 후 5건 일괄 재시험(Pass 전환 유력).

\textbf{유형 B(4건).} TC-BLE-001은 매칭/비매칭 BLE 기기가 동시에 다수 필요한 필터 배제
검증으로 테스트 베드 미구성, 아울러 부팅 시 \code{\_sync\_bt\_name()}이 광고 이름을
``ARES''로 재설정해 시점(최초/재부팅)에 따라 이름/UUID 매칭이 달라지는 구조적 혼동
지점도 확인. TC-GUN-001$\sim$003은 총 모듈(GP22 PWM) 미장착으로 물리 검증 미수행 ---
다만 시뮬레이터 발사 애니메이션(TC-SIM-016$\cdot$023)과 fire-and-forget 일반 규칙
(TC-BLE-004), STOP\_ALL 전 하드웨어 정지(TC-CONN-007)는 Pass로 경로 골격은 검증 완료.

\textbf{유형 C(2건).} TC-BLE-006(FIRE\_AND\_FORGET 양쪽 동기화)은 소스코드 정적 대조
케이스인데 TCL에 대조 파일$\cdot$상수명$\cdot$방법이 명시되지 않아 수행 불가(``이해를
못했습니다'') --- 테스트 설계와 테스터 역량 매칭 실패 사례. v1.0 코드 분석에서는 28개
head 일치를 확인한 바 있으나 독립 검증 미완으로 보류 유지. TC-BLK-005는 참조 프리셋
명칭 혼동으로 기타 처리(기능 자체는 ``리턴 값 없는 함수 실행 작동'' 확인, 사실상 Pass에 준함).

\textbf{유형 D(2건).} \code{LED\_PATTERN}(swipe\_effect)$\cdot$\code{SET\_PIN}은 펌웨어에
구현되어 있으나 블록/Web UI에 실행 경로가 노출되지 않은 \emph{데드 코드/관리자 기능} ---
학습자 시나리오 영향은 없으며, 블록 신설$\cdot$명령 제거$\cdot$UI 노출 여부의 제품 결정
(IMP-002$\cdot$003) 후 케이스를 재정의한다.

\section{성능 KPI · 결함/개선 항목}
KPI 5개 항목은 계측 수단 제약(rx 로그 부재, 오실로스코프 미투입)으로 정밀 실측을 차기
회차로 이월했다. 다만 \textbf{정상 명령 응답 timeout 0건}은 실행 전 케이스에서 관찰
기준으로 충족되었고, DC 양방향 대칭은 육안상 뚜렷한 비대칭 없음(TC-DCM-002 Pass),
BLE 페어링은 체감상 즉시 연결 수준으로 소견 기록되었다.

\begin{table}[htbp]\centering\small
\caption{결함(BUG) 및 개선 항목(IMP)}
\renewcommand{\arraystretch}{1.2}
\begin{tabularx}{\linewidth}{@{}l l X l@{}}
\toprule
\sffamily ID & \sffamily 관련 & \sffamily 내용 & \sffamily 우선순위 \\
\midrule
BUG-001 & TC-BLE-008 & UART 9600 baud 고정 --- 처리량의 구조적 천장(차기 이관, 38400+ 상향 검토) & Mid \\
BUG-002 & --- & 배터리 전압 감시 미구현(\code{get\_battery\_voltage()} 상수 100) --- 차기 이관 & Low \\
\midrule
IMP-001 & 보류 5건 & 실기 BLE 경로의 \textbf{수신(rx) 문자열 로그 표시} + 연결 수립 단계 로그 & \textbf{High} \\
IMP-002 & TC-LED-005 & \code{LED\_PATTERN} 레거시 정리(블록 신설 또는 명령 제거 택일) & Low \\
IMP-003 & TC-SYS-004 & \code{SET\_PIN} Web 노출 여부 결정(설정 UI 또는 관리자 기능 확정) & Low \\
IMP-004 & C 유형 등 & TCL 절차 보강(소스 대조 방법$\cdot$BLE 필터 시점별 판정$\cdot$프리셋 경로 명시) & Mid \\
\bottomrule
\end{tabularx}
\end{table}

신규 발견 기능 결함(Fail)은 \textbf{0건}이며(심각도 전 구간 0), 잔여 Mid/Low 2건은 모두
사전 식별된 구조적 한계로 차기 버전 이관을 유지한다. 중점 점검(부팅 None 가드, 비상 정지
ABORT, BATCH 금지 블록 이중 방어, 설정 영구 저장, BLE 끊김 안전 중지, 시뮬레이터 전
경로)은 전건 Pass.

\section{재시험 계획 및 최종 판정}
\begin{table}[htbp]\centering\small
\caption{재시험 계획}
\renewcommand{\arraystretch}{1.2}
\begin{tabularx}{\linewidth}{@{}l X X@{}}
\toprule
\sffamily 차수 & \sffamily 대상 & \sffamily 선결 조건 \\
\midrule
R1 & 보류 5건(유형 A) + KPI 1$\cdot$2$\cdot$4 실측 & IMP-001(rx 로그) 반영 빌드 \\
R2 & TC-BLE-006, TC-BLK-005 & IMP-004(TCL 절차 보강), 개발 담당 배석 \\
R3 & TC-GUN-001$\sim$003 & 총 모듈 수급(안전 케이스 TC-GUN-003 우선) \\
R4 & TC-BLE-001 + KPI 3 정밀 & BLE 기기 3종 이상 테스트 베드, 계측 장비 \\
\bottomrule
\end{tabularx}
\end{table}

\noindent\textbf{최종 판정 --- 조건부 Go.} 실행 85건 전건 Pass$\cdot$Fail 0건, 하드웨어
구동$\cdot$시뮬레이터 전 경로$\cdot$E2E 정상 완주, 안전$\cdot$안정성 항목 전건 통과로
기능 품질은 배포 수준으로 평가한다. 다만 High 보류 5건의 독립 검증이 남아 있어
\textbf{재시험 R1(rx 로그 반영)$\cdot$R2(절차 보강) 완료를 조건으로 Go 를 확정}하며,
R3$\cdot$R4는 현장 실증 단계와 병행을 허용한다. (승인: 품질관리 책임자)